\section{Conditioned kernels and quantitative Gaussian comparison}
\label{sec:quantitative-comparison}

Fix the input dimension $d$, depth $L$, and margin $a>0$. Put
$M=2^d$, $b_N=N^{(1-\gamma)/2}$ and $t_N=N^\gamma$, and let $Q_N$
be the empirical covariance matrix of the last hidden layer. All norms
in this section are the maximum entry norm, unless another norm is
indicated. Constants may depend on $d,L,a$; none of the statements is
uniform in these parameters. The exact independent Gaussian readout gives
\[
 Z_N\mid Q_N\sim\mathcal N(0,Q_N/t_N).
\]
For a full Boolean labeling $g$, set
\[
 O_g=\{z:g(x)z(x)>a\text{ for every }x\}.
\]
For a collection $S$ of observed labels, let $\mathcal G_S$ be the
nonempty set of full labelings consistent with $S$, and let
$A_S=\bigcup_{g\in\mathcal G_S}O_g$. Thus conditioning on $A_S$
includes the margin requirement on every unobserved input.

\subsection{The finite-gain center and conditional concentration}

For a covariance matrix $R$, write
\[
 \Phi_b(R)_{xy}=\mathbb E[\tanh(bU_x)\tanh(bU_y)],
 \qquad U\sim\mathcal N(0,R).
\]
Define deterministic matrices
\[
 R_0(x,y)=\frac{1+x\cdot y}{d},\qquad
 K_N^1=\Phi_{b_N}(R_0),\qquad
 K_N^s=\Phi_{b_N}(K_N^{s-1}),\qquad
 \widehat K_N=K_N^{L-1}.
\]
The distinction between $\widehat K_N$ and $\mathbb E Q_N$ matters:
at depth greater than two the nonlinear covariance recursion does not
commute with expectation. The limiting saturated kernel is
\[
 K^{(0)}(x,y)=\frac{1+x\cdot y}{d+1},\qquad
 K^{(s+1)}=\frac2\pi\arcsin K^{(s)},\qquad K=K^{(L-1)},
\]
where the arcsine is applied entrywise. This matrix is strictly positive
definite. Indeed, the first-layer sign features associated to the
nonparallel augmented inputs $(1,x)$ are linearly independent: crossing
one of their distinct hyperplanes changes only its corresponding sign.
Their Gram matrix is therefore positive definite. Subsequent arcsine
maps preserve this property, by their positive power series and the
Schur product theorem.

\begin{lemma}[Concentration about the finite-gain kernel]
\label{lem:cmp-concentration}
For each fixed $\gamma<1$, there are positive constants $C,c,u_0$ and
$N_0$ such that, for $N\ge N_0$ and $0<u\le u_0$,
\begin{align}
 \Pr\{\|Q_N-\widehat K_N\|>u\}&\le C e^{-cNu^2},
 \label{eq:cmp-finite-concentration}\\
 \|\widehat K_N-K\|&\le C b_N^{-1}.
 \label{eq:cmp-finite-bias}
\end{align}
\end{lemma}

\begin{proof}
The diagonal variances of all limiting hidden kernels are one, and
their distinct-coordinate correlations have absolute value strictly
less than one. In fixed neighborhoods of these matrices, $\Phi_b$ has
a Lipschitz constant independent of $b$. To verify the uniformity, each
off-diagonal entry uses only a two-dimensional Gaussian marginal.
Differentiate its density with respect to its covariance parameters.
On a compact set of positive variances and correlations bounded away
from $\pm1$, the $L^1$ norms of these density derivatives are bounded.
The integrand $\tanh(bu)\tanh(bv)$ has absolute value at most one,
so the resulting derivative bound is independent of $b$. The same
one-dimensional argument handles diagonal entries. This argument does
not require the full input covariance $R_0$ to be invertible.

Given the preceding hidden layer, its $N$ successor rows are independent.
Each entry of their empirical covariance averages $N$ variables in
$[-1,1]$. Conditional Hoeffding bounds, a union bound over the $M^2$
entries, and the uniform Lipschitz estimate give
\eqref{eq:cmp-finite-concentration} by induction over the fixed number
of layers. At each stage one applies the previous-layer estimate with
a constant multiple of $u$; exits from the fixed neighborhoods have
probability at most $C e^{-cN}$ and are absorbed into the same bound.

For a centered Gaussian $U$ of variance bounded below,
\[
 \mathbb E|\tanh(bU)-\operatorname{sign}U|\le C/b.
\]
The Gaussian sign-correlation formula and the same Lipschitz induction
therefore give \eqref{eq:cmp-finite-bias}.
\end{proof}

\begin{proposition}[Conditioning preserves the frozen kernel]
\label{prop:cmp-conditional-freezing}
For every fixed $\gamma<1$,
\[
 Q_N\longrightarrow K
 \quad\text{in probability conditional on }A_S.
\]
For $0<\gamma<1$, there is $C_0$ such that
$\Pr(A_S)\ge e^{-C_0t_N}$ for all sufficiently large $N$.
Consequently, for fixed $\epsilon>0$,
\[
 \Pr\{\|Q_N-K\|>\epsilon\mid A_S\}
 \le C_\epsilon e^{-c_\epsilon N+C_0t_N}.
\]
More quantitatively, for a fixed sufficiently large constant $A$,
\[
 \Pr\left\{\|Q_N-K\|>
 C b_N^{-1}+A\sqrt{t_N/N}\,\middle|\,A_S\right\}
 \le C\exp[-(cA^2-C_0)t_N].
\]
These bounds may be chosen simultaneously for all observed-label
collections $S$ at the fixed input dimension.
\end{proposition}

\begin{proof}
Choose any $g\in\mathcal G_S$ and any fixed bounded box strictly inside
$O_g$. On a sufficiently small fixed neighborhood of $K$, all covariance
eigenvalues are bounded above and below. Integrating the density of
$\mathcal N(0,Q_N/t_N)$ over the box gives a lower bound
$c t_N^{M/2}e^{-Ct_N}$. That kernel neighborhood has probability tending
to one, proving the asserted denominator bound. Dividing the bounds of
Lemma~\ref{lem:cmp-concentration} by this denominator proves the claims
for positive $\gamma$. Here $t_N=o(N)$, and
$b_N^{-1}=\sqrt{t_N/N}$ tends to zero.

If $\gamma=0$, the normalized output converges to the nondegenerate
Gaussian $\mathcal N(0,K)$ and the conditioning event has a positive
limiting probability. If $\gamma<0$, its normalized margin is
$a\sqrt{t_N}\to0$; the limiting conditioning event requires only the
observed signs, again with positive Gaussian probability. The
unconditional kernel concentration proves the corresponding conditional
claim in both cases.
\end{proof}

\subsection{What holds at every positive subcritical exponent}

Set
\[
 I_K(z)=\frac12z^TK^{-1}z,\qquad
 E_K(g)=\min_{z\in\overline O_g}I_K(z),\qquad
 E_S=\min_{g\in\mathcal G_S}E_K(g).
\]
Each orthant has a unique minimizing vector $z_g$. Write
\[
 \mathcal H_S=\{g\in\mathcal G_S:E_K(g)=E_S\},\qquad
 V_S=\{z_g:g\in\mathcal H_S\}.
\]

\begin{proposition}[Conditional latent concentration and unique minima]
\label{prop:cmp-minimizers}
For every $0<\gamma<1$, the posterior concentrates on $V_S$. More
precisely, for every $\delta>0$ there is $\eta_\delta>0$ such that
\[
 \limsup_{N\to\infty}t_N^{-1}\log
 \Pr\{\operatorname{dist}(Z_N,V_S)\ge\delta\mid A_S\}
 \le-\eta_\delta.
\]
If $\mathcal H_S=\{g_*\}$, the posterior on full Boolean functions
converges in total variation to $\delta_{g_*}$. The fixed Gaussian
$\mathcal N(0,K/t_N)$, conditioned on the same $A_S$, has the same
limit. Thus their Boolean posteriors agree asymptotically throughout
$0<\gamma<1$ whenever the compatible energy minimum is unique.
\end{proposition}

\begin{proof}
Couple the output to $Y_N=K^{1/2}\xi/\sqrt{t_N}$ by writing
$Z_N=Q_N^{1/2}\xi/\sqrt{t_N}$, with $\xi$ a standard independent
$M$-dimensional Gaussian. Near the positive definite $K$, the matrix
square root is Lipschitz. For fixed $\epsilon,\delta>0$, kernel
concentration and the Gaussian tail therefore bound
\[
 \Pr\{\|Z_N-Y_N\|_2>\delta\}
 \le C_\epsilon e^{-c_\epsilon N}
 +\Pr\{\|\xi\|_2>\delta\sqrt{t_N}/(C\epsilon)\}.
\]
After division by $t_N$ in the logarithm, let $N\to\infty$ and then
$\epsilon\downarrow0$. Since $t_N=o(N)$, this proves exponential
equivalence at speed $t_N$. Hence $Z_N$ has the Gaussian large-deviation
rate $I_K$.

The minimum rate on the strict orthant equals that on its closure:
multiply its minimizing vector by $1+\varepsilon$ and let
$\varepsilon\downarrow0$. The full-margin denominator consequently
has logarithmic rate $E_S$. Coercivity of $I_K$ and the finiteness of
$\mathcal G_S$ give a strictly positive rate gap on the closed feasible
set at distance at least $\delta$ from $V_S$. The large-deviation upper
bound and the matching denominator lower bound give the displayed
conditional estimate. If the minimizing labeling is unique, the same
argument gives an exponential gap to every other compatible labeling.
Apply the identical argument to $Y_N$.
\end{proof}

If several labelings minimize the energy, this proposition fixes their
possible support but leaves their relative posterior weights undetermined.
Those weights can decide whether a correct label has probability greater
than $2/3$. The next result controls them on a smaller exponent range.

\subsection{Relative probabilities and total variation}

\begin{lemma}[Covariance sensitivity of an orthant probability]
\label{lem:cmp-log-lipschitz}
Let $p_g(Q,t)=\Pr\{\mathcal N(0,Q/t)\in O_g\}$. In a sufficiently
small fixed convex neighborhood of $K$, uniformly over $g$ and $t\ge1$,
\[
 |\log p_g(Q,t)-\log p_g(Q',t)|\le Ct\|Q-Q'\|.
\]
\end{lemma}

\begin{proof}
Uniformly over this covariance neighborhood, integrating over a fixed
box inside each orthant gives $p_g(Q,t)\ge e^{-C_1t}$, with constants
chosen for all finitely many $g$. An unconditional Gaussian tail bound
gives
\[
 \Pr\{\|Z\|_2^2>u\}\le C e^{-ctu}
\]
for $u$ above a fixed constant. Divide by the preceding lower bound
and integrate the resulting conditional tail, starting at a sufficiently
large fixed $u$. It follows that
$\mathbb E[\|Z\|_2^2\mid Z\in O_g]\le C$ uniformly for $t\ge1$.
Differentiating the Gaussian orthant integral along the segment from
$Q$ to $Q'$ now gives a determinant derivative bounded by
$C\|Q-Q'\|$ and a conditional quadratic-moment derivative bounded by
$Ct\|Q-Q'\|$. Integration along the segment proves the claim.
\end{proof}

\begin{theorem}[Quantitative Gaussian comparison]
\label{thm:cmp-relative}
For $0<\gamma<1$, put $t=t_N$. There are positive constants $C,c$
such that
\[
 \max_g\left|
 \frac{\Pr\{Z_N\in O_g\}}{p_g(\widehat K_N,t)}-1\right|
 \le C\frac{t}{\sqrt N}\exp(Ct^2/N)+C\exp(-cN+Ct)
\]
for all sufficiently large $N$. In particular, if $0<\gamma<1/2$,
the Boolean posterior of $Z_N$ conditional on $A_S$ is $o(1)$ in total
variation from that of
$\mathcal N(0,\widehat K_N/t_N)$ conditional on $A_S$.
If $0<\gamma<1/3$, the same statement holds with $K$ in place of
$\widehat K_N$. These assertions include tied minimum-energy labelings.
\end{theorem}

\begin{proof}
Put $D=\|Q_N-\widehat K_N\|$ and restrict first to $D\le u_0$,
with $u_0$ chosen small enough for
Lemma~\ref{lem:cmp-log-lipschitz}. On that event,
\[
 \left|\frac{p_g(Q_N,t)}{p_g(\widehat K_N,t)}-1\right|
 \le e^{CtD}-1.
\]
Integration of the sub-Gaussian tail
\eqref{eq:cmp-finite-concentration} yields
\[
 \mathbb E[(e^{CtD}-1)\mathbf 1_{D\le u_0}]
 \le C\frac{t}{\sqrt N}\exp(Ct^2/N).
\]
For example, integrate
$Ct\int_0^{u_0}e^{Ctu}\Pr(D>u)\,du$ and complete the square;
enlarging the integral to $[0,\infty)$ gives the displayed bound.
On $D>u_0$, the numerator is at most one, its probability is
$Ce^{-cN}$, and the reference orthant probability is at least
$e^{-Ct}$. This proves the relative-probability estimate, including
the negligible subtraction needed in its lower bound.

When $t=o(\sqrt N)$, the estimate tends to zero uniformly over all
full sign patterns. Summing over $\mathcal G_S$ and dividing by the
corresponding sums shows that conditional Boolean probabilities agree
in total variation, uniformly over $S$. Finally,
Lemmas~\ref{lem:cmp-concentration} and
\ref{lem:cmp-log-lipschitz} show that replacing $\widehat K_N$ by
$K$ changes each orthant probability by a factor
$\exp[O(t/b_N)]$. This factor tends to one when
$t/b_N=N^{(3\gamma-1)/2}\to0$, proving the last assertion.
\end{proof}

The two comparisons address distinct errors. Empirical fluctuations of
the hidden kernel lead to the sufficient condition $\gamma<1/2$.
Replacing the deterministic finite-gain kernel by its saturated limit
introduces a separate bias, for which $\gamma<1/3$ is a sufficient
condition. Conditional kernel convergence and the leading Gaussian
large-deviation rate remain valid over the entire range $\gamma<1$.
